\documentclass[11pt,a4paper]{article}

\usepackage[T1]{fontenc}
\usepackage[utf8]{inputenc}
\usepackage{amsmath}
\usepackage{booktabs}
\usepackage{enumitem}
\usepackage{geometry}
\usepackage{hyperref}
\usepackage{xcolor}
\usepackage{listings}
\usepackage{tabularx}

\geometry{margin=1in}
\setlength{\emergencystretch}{3em}
\setlist{nosep}
\hypersetup{
    colorlinks=true,
    linkcolor=blue!50!black,
    urlcolor=blue!50!black
}

\lstset{
    basicstyle=\ttfamily\small,
    columns=fullflexible,
    keepspaces=true,
    frame=single,
    breaklines=true,
    backgroundcolor=\color{black!3}
}

\newcommand{\code}[1]{\texttt{\detokenize{#1}}}
\newcommand{\ms}[1]{\ensuremath{#1\,\text{ms}}}

\title{GPU Performance Report}
\author{Henrique Rocha}
\date{\today}

\begin{document}
\maketitle

\begin{abstract}
This document analyses GPU-side performance of the Vulkan rendering engine at the time of writing. It describes the frame breakdown, identifies bottlenecks in each sub-renderer, quantifies estimated improvement for each issue, and provides a prioritized action plan. A GPU timestamp query infrastructure already exists and reports per-pass timings for shadow, cull, sky, solid draw, vegetation, water, and ImGui.
\end{abstract}

\tableofcontents

\section{Engine Overview}

The engine renders terrain, vegetation, water, and sky into offscreen targets and composits them via a post-process pass. It uses dynamic rendering (\code{VK\_KHR\_dynamic\_rendering}) throughout, GPU-driven frustum culling via compute shaders, and cascaded shadow maps.

\subsection{Render Graph (per frame)}

\begin{enumerate}
    \item GPU culling dispatches (solid geometry, vegetation)
    \item Shadow pass: 3 cascades $\times$ depth-only (solid + vegetation)
    \item Solid offscreen: optional depth pre-pass, color + depth pass
    \item Vegetation rendering: depth pre-pass, color pass, impostor pass
    \item Water geometry pass into offscreen target
    \item Water back-face depth pass (async)
    \item 360\textdegree{} cubemap capture (async)
    \item Post-process composite into swapchain
\end{enumerate}

\subsection{Key Constants}

\begin{center}
\begin{tabular}{llr}
\toprule
\textbf{Parameter} & \textbf{Value} \\
\midrule
Shadow cascades & 3 \\
Shadow map resolution & $2048 \times 2048$ \\
Offscreen solid format & Swapchain format ($32$ bpp) \\
Offscreen depth format & \code{VK\_FORMAT\_D32\_SFLOAT} \\
Water offscreen format & \code{VK\_FORMAT\_R32G32B32A32\_SFLOAT} ($128$ bpp) \\
Frames in flight & 2 \\
Cull overlap frames & 3 \\
Texture array layers & 32 (max) \\
Texture array resolution & $1024 \times 1024$ \\
\bottomrule
\end{tabular}
\end{center}

\section{GPU Timestamp Infrastructure}

The application records seven GPU intervals per frame via \code{VK\_QUERY\_TYPE\_TIMESTAMP} (see \code{main.cpp} lines 81--95):

\begin{center}
\begin{tabular}{cl}
\toprule
\textbf{Interval} & \textbf{GPU work covered} \\
\midrule
0--1 & Shadow pass (all cascades) \\
2--3 & GPU culling dispatches \\
4--5 & Sky rendering \\
6--7 & Solid draw \\
8--9 & Vegetation draw \\
10--11 & Water draw \\
12--13 & ImGui rendering \\
\bottomrule
\end{tabular}
\end{center}

Results are read back via \code{vkGetQueryPoolResults} once per frame and displayed in the settings UI. The \code{timestampPeriod} from \code{VkPhysicalDeviceProperties} is used to convert to milliseconds.

\section{Performance Issues}

\subsection{Shadow Mapping}

\paragraph{5$\times$5 PCF on 3 Cascades (Critical)}

\textbf{Files:} \code{shaders/includes/shadows.glsl:2--13}, \code{shaders/main.frag:147--156}

The terrain fragment shader uses a $5 \times 5$ percentage-closer filter (25 taps) per cascade. A fragment that cascades through all three shadow maps issues up to 75 texture samples for shadow alone. The shadow maps use \code{VK\_FILTER\_NEAREST}, so no hardware bilinear filtering assists.

\textbf{Solution:} Switch to variance shadow maps (VSM), or reduce PCF to $3 \times 3$ with linear filtering, or use a single-sample hard shadow for distant fragments.

\textbf{Estimated improvement:} $3{-}8\times$ shadow shader throughput.

\paragraph{No Re-Culling Per Cascade (Medium)}

\textbf{Files:} \code{vulkan/renderer/SceneRenderer.cpp:shadowPass()}

Solid geometry in the shadow pass reuses the indirect buffer that was culled against the camera frustum. All visible meshes are rendered into all three cascades, even when a mesh lies entirely outside a given cascade's light frustum.

\textbf{Solution:} Submit three culling dispatches, one per cascade matrix, and render each cascade from its own compacted indirect buffer.

\textbf{Estimated improvement:} $1{-}3\ms$ for the shadow pass on complex scenes.

\paragraph{Depth Bias Pipeline Flag With Zero Bias (Low)}

\textbf{Files:} \code{vulkan/renderer/ShadowRenderer.cpp}

The shadow pipeline is created with \code{depthBiasEnable = true}, but \code{vkCmdSetDepthBias} is called with $(0.0, 0.0, 0.0)$ every frame. The flag triggers a state change for no benefit.

\textbf{Solution:} Create the pipeline with \code{depthBiasEnable = false} and remove the \code{vkCmdSetDepthBias} call.

\subsection{Terrain Rendering (Solid Renderer)}

\paragraph{Tessellation in Shadow and Depth Passes (Medium)}

\textbf{Files:} \code{shaders/main.tesc}, \code{shaders/main.tese}, \code{vulkan/renderer/SolidRenderer.cpp}

The tessellation control and evaluation shaders run for every vertex in the shadow pass and depth pre-pass, evaluating displacement mapping even though only depth is needed.

\textbf{Solution:} Create a no-tessellation depth-only pipeline that uses a simple vertex shader and skips TCS/TES entirely.

\textbf{Estimated improvement:} $1{-}2\ms$ combined for shadow and depth passes.

\subsection{Vegetation Rendering}

\paragraph{Geometry Shader for Billboard Expansion (High)}

\textbf{Files:} \code{shaders/vegetation.geom}, \code{vulkan/renderer/VegetationRenderer.cpp}

Billboard quads are expanded from point primitives in a geometry shader. Geometry shaders perform poorly on modern GPU architectures (NVIDIA Turing/Ada, AMD RDNA) compared to vertex-shader-based expansion.

\textbf{Solution:} Expand billboard quads in the vertex shader by passing pre-computed corner offsets, or use a compute shader to generate screen-aligned quads.

\textbf{Estimated improvement:} $2{-}4\times$ vegetation shading throughput.

\paragraph{Double Rasterization (VK\_CULL\_MODE\_NONE) (High)}

\textbf{Files:} \code{vulkan/renderer/VegetationRenderer.cpp}

All vegetation pipelines use \code{VK\_CULL\_MODE\_NONE}, rasterizing both front and back faces. For camera-facing billboard quads, approximately 50\% of fragments are back-facing and discarded only after the fragment shader.

\textbf{Solution:} Use \code{VK\_CULL\_MODE\_BACK\_BIT}.

\textbf{Estimated improvement:} $1.5{-}2\times$ vegetation fill rate.

\paragraph{Wind Push Constants Exceed 128 Bytes (Low)}

\textbf{Files:} \code{shaders/vegetation.vert}, \code{shaders/impostors.frag}

Wind parameters passed as push constants are approximately 160 bytes. The Vulkan minimum guaranteed limit is 128 bytes. This works on desktop GPUs but may fail on mobile or integrated hardware.

\textbf{Solution:} Move wind parameters to a UBO. Not a bottleneck on current target hardware.

\subsection{Water Rendering}

\paragraph{128 bpp Offscreen Target (High)}

\textbf{Files:} \code{vulkan/renderer/WaterRenderer.cpp}

The water color target is \code{VK\_FORMAT\_R32G32B32A32\_SFLOAT} (16 bytes per pixel). Only the alpha channel carries linear depth. This is four times the bandwidth of a typical R8G8B8A8\_UNORM + separate R32\_SFLOAT depth combination.

\textbf{Solution:} Use \code{VK\_FORMAT\_R16G16B16A16\_SFLOAT} (8 bytes/pixel) or \code{VK\_FORMAT\_R8G8B8A8\_UNORM} with a separate R32 depth attachment.

\textbf{Estimated improvement:} $2{-}4\times$ water memory bandwidth, approximately $0.5{-}1\ms$ saved.

\paragraph{Full-Screen Depth Copy (Medium)}

\textbf{Files:} \code{vulkan/renderer/WaterRenderer.cpp:initializeGeomDepthFromSceneDepth()}

The scene depth buffer is copied to the water geometry depth buffer via \code{vkCmdCopyImage}. This is a full-resolution bandwidth operation.

\textbf{Solution:} Use the same depth attachment for water rendering with \code{VK\_ATTACHMENT\_LOAD\_OP\_LOAD} after transitioning the solid depth image to the correct layout.

\textbf{Estimated improvement:} $0.5{-}1\ms$ per frame.

\paragraph{Cubemap Water Without Culling (Low)}

\textbf{Files:} \code{vulkan/renderer/WaterRenderer.cpp}

When rendering water into the 360\textdegree{} cubemap for reflections, \code{drawAll} is used (no frustum culling). This draws all water patches six times.

\textbf{Solution:} Cull water patches per cubemap face using the face's view-projection matrix.

\subsection{GPU Culling (IndirectRenderer)}

\paragraph{Host-Visible Cull Count Buffers (Low--Medium)}

\textbf{Files:} \code{vulkan/renderer/IndirectRenderer.hpp:153--155}

Visible count buffers are allocated as \code{HOST\_VISIBLE | HOST\_COHERENT} and reset via a CPU pointer write inside \code{prepareCull}. On discrete GPUs this requires uncached PCIe reads.

\textbf{Solution:} Use \code{vkCmdFillBuffer} or a dedicated compute shader to reset the counter.

\textbf{Estimated improvement:} $0.2{-}0.5\ms$ per frame.

\subsection{Synchronization and Uploads}

\paragraph{queueWaitIdle() During Mesh Uploads (Medium)}

\textbf{Files:} \code{vulkan/renderer/IndirectRenderer.cpp}

The \code{rebuild()} and \code{uploadMeshVerticesAndIndices()} paths call \code{queueWaitIdle()} before staging copies, stalling the entire graphics pipeline.

\textbf{Solution:} Use fence-based deferred publication: submit transfers, return a fence, and defer buffer-pointer publication until the fence signals.

\textbf{Estimated improvement:} Removes frame drops during scene loading and dynamic mesh updates.

\paragraph{UBO Map/Unmap Every Frame (Low)}

\textbf{Files:} \code{main.cpp:715--722}

The per-frame UBO is updated by mapping, \code{memcpy}, then unmapping the buffer. This causes a validation-layer map/unmap pair each frame.

\textbf{Solution:} Persistently map the UBO at allocation time. Use a ring of mapped pointers or \code{vkMapMemory} once during init.

\section{Aggregate Improvement Estimate}

The following table estimates the combined effect on a mid-range GPU (RTX 3060 / RX 6600) for a complex outdoor scene.

\begin{center}
\begin{tabular}{lcc}
\toprule
\textbf{Pass} & \textbf{Current (est.)} & \textbf{Optimised (est.)} \\
\midrule
Shadow pass (3 cascades) & $4{-}6\ms$ & $1.5{-}2\ms$ \\
GPU culling & $0.3{-}0.5\ms$ & $0.2{-}0.3\ms$ \\
Sky & $0.2{-}0.4\ms$ & $0.2{-}0.4\ms$ \\
Solid draw & $2{-}4\ms$ & $1.5{-}3\ms$ \\
Vegetation & $3{-}5\ms$ & $1{-}2\ms$ \\
Water & $2{-}4\ms$ & $1{-}2\ms$ \\
ImGui & $0.5{-}1\ms$ & $0.5{-}1\ms$ \\
\midrule
\textbf{Total GPU} & $12{-}18\ms$ & $6{-}10\ms$ \\
\textbf{Frame rate} & 55--83 FPS & 100--166 FPS \\
\bottomrule
\end{tabular}
\end{center}

\section{Recommended Action Plan}

\begin{enumerate}
    \item \textbf{Switch shadow PCF to $3\times3$ with linear filtering, or adopt VSM.} This is the single largest bottleneck and touches every fragment that receives shadows.
    \item \textbf{Eliminate the geometry shader for vegetation billboards.} Move expansion to the vertex shader. This alone doubles vegetation throughput.
    \item \textbf{Set \code{VK\_CULL\_MODE\_BACK\_BIT} on all vegetation pipelines.} A one-line change with immediate fill-rate benefit.
    \item \textbf{Reduce water offscreen format to R16G16B16A16\_SFLOAT.} Cuts water memory bandwidth by $2{-}4\times$.
    \item \textbf{Re-cull solids per shadow cascade.} Submit three compute dispatches with cascade-specific matrices.
    \item \textbf{Share depth buffer with water instead of copying.} Eliminates a full-screen \code{vkCmdCopyImage}.
    \item \textbf{Replace \code{queueWaitIdle()} with fence-based deferred publication.} Prevents GPU bubbles during mesh uploads.
    \item \textbf{Create a no-tessellation depth-only pipeline for shadow and depth passes.}
    \item \textbf{Reset cull count via \code{vkCmdFillBuffer} instead of host-mapped write.}
    \item \textbf{Generate mipmaps for texture arrays.} Improves texture cache behaviour when sampling at distance.
\end{enumerate}

\section{Shader Complexity Summary}

\begin{center}
\begin{tabular}{lcc}
\toprule
\textbf{Shader} & \textbf{Lines} & \textbf{Notes} \\
\midrule
\code{main.frag} & 498 & Triplanar mapping, $5\times5$ PCF shadows, TBN, env reflection \\
\code{main.vert} & 49 & Standard pass-through \\
\code{main.tesc} & 114 & Per-edge tess factor, texture index compression \\
\code{main.tese} & 97 & Displacement mapping, sharp normal computation \\
\code{water.frag} & 605 & Perlin noise, fresnel, caustics, reflection, refraction \\
\code{vegetation.frag} & 130 & Billboard alpha, normal mapping, impostor fade \\
\code{sky.frag} & 74 & Procedural gradient, sun flare, stars \\
\code{postprocess.frag} & 75 & Sky--solid--water composite \\
\code{indirect.comp} & 117 | 89 & GPU frustum culling compute shaders \\
\bottomrule
\end{tabular}
\end{center}

The \code{water.frag} shader at 605 lines is the longest in the project and represents a significant GPU workload, especially given the 128 bpp offscreen target it writes to.

\end{document}
